% Abstract, shared by main.tex (TCBB) and biorxiv.tex. About 200 words.
The Gene Expression Omnibus holds sequencing data from more than three billion cells, and machine learning models now
train on millions of them at a time, streamed from cloud storage in minibatches. For these workflows the bottleneck is
moving and decoding count matrices, not computing on them. We measured how few bits these counts need and how fast they
can be decoded, then built \sone{} to approach both limits, using Claude 5.5-generation models (Anthropic). A \sone{}
dataset is a manifest of content-addressed sample files, each opening with an index to blocks of cells, so a reader
fetches only the cells it needs by HTTP range request. Within a block, genes are numbered by detection frequency, so each
cell becomes small steps and small counts, which an rANS coder encodes with tables chosen by the cell's depth and
position along the gene list, in 32 lanes for SIMD decoding. On \nm{nMatrices} public matrices, \sone{} files were the
smallest on every one, \nm{ratioBestOther} times smaller than the best alternative and within a few percent of
cross-validated count models; in UMI data most remaining bits are molecular sampling, irreducible given each cell's
expression rates. On four cores \sone{} read fastest on \nm{dec-s1-fastest-count}, and from a real cloud object store it
delivered a first PyTorch minibatch in under a second and streamed cells tens of times faster than Zarr, though per core
its decoder is slower than byte-oriented codes.
